% XRPL Lending Protocol Hackathon · CY-HACK · Developer Feedback (3 pages)
% Build : tectonic -X compile bonus/deck-source/feedback-report.tex --outdir bonus/deck-source && mv bonus/deck-source/feedback-report.pdf FEEDBACK.pdf
% Captures : node bonus/deck-source/capture-explorer.mjs   (→ bonus/deck-source/shots/*.png)
% DA du deck : navy 001C5C · bleu 006AFF · ciel 6DC3FF · corps 5F666E
%
% Chaque constat porte les champs de la consigne : title, category, severity, library
% and version, description, repro steps, transaction/code links, proposed fix.
% Règles de rédaction : une information à un seul endroit ; la description dit le
% problème, la repro donne les valeurs mesurées ; une seule icône, la sévérité.
\documentclass[a4paper,10pt]{extarticle}

\usepackage[left=13mm,right=13mm,top=11mm,bottom=13mm,footskip=6mm]{geometry}
\usepackage{fontspec}
\usepackage[table]{xcolor}
\usepackage{tabularx,booktabs,array,colortbl}
\usepackage{multicol}
\usepackage{graphicx}
\usepackage{xfp}
\usepackage{xstring}
\usepackage{tcolorbox}
\usepackage{enumitem}
\usepackage{microtype}
\usepackage{fancyhdr}
\usepackage{ragged2e}
\usepackage{lastpage}
\usepackage{fontawesome5}
\usepackage[hidelinks]{hyperref}
\hypersetup{pdftitle={XRPL Lending Protocol: Developer Feedback, CY-HACK},pdfauthor={CY-HACK}}
\tcbuselibrary{skins}
\graphicspath{{bonus/deck-source/shots/}{shots/}}

% ── DA ──
\definecolor{navy}{HTML}{001C5C}
\definecolor{blue}{HTML}{006AFF}
\definecolor{sky}{HTML}{6DC3FF}
\definecolor{body}{HTML}{5F666E}
\definecolor{muted}{HTML}{8A9096}
\definecolor{card}{HTML}{F6F7F8}
\definecolor{line}{HTML}{E2E5E8}
\definecolor{bluecard}{HTML}{EFF6FF}
% sévérité : texte, fond de pastille, filet gauche de la carte
\definecolor{highfg}{HTML}{B93A1F}\definecolor{highbg}{HTML}{FBE9E4}\definecolor{highrule}{HTML}{E0532F}
\definecolor{medfg}{HTML}{8A5700}\definecolor{medbg}{HTML}{FDF1D6}\definecolor{medrule}{HTML}{E2A21B}
\definecolor{lowfg}{HTML}{4F5B66}\definecolor{lowbg}{HTML}{ECEFF2}\definecolor{lowrule}{HTML}{A7B0B8}

% Faces déclarées explicitement : laissée à la résolution automatique,
% HelveticaNeue.ttc sort \bfseries en Bold Italic. Pas d'option Color= non plus :
% elle casse xdvipdfmx avec tcolorbox (« Invalid object type »).
\setmainfont{Helvetica Neue}[
  UprightFont = {Helvetica Neue}, BoldFont = {Helvetica Neue Bold},
  ItalicFont = {Helvetica Neue Italic}, BoldItalicFont = {Helvetica Neue Bold Italic}]
\setmonofont{Menlo}[Scale=0.84,
  UprightFont = {Menlo Regular}, BoldFont = {Menlo Bold},
  ItalicFont = {Menlo Italic}, BoldItalicFont = {Menlo Bold Italic}]
\newfontfamily\hdfont{Poppins}[BoldFont={Poppins SemiBold}]
\makeatletter\renewcommand\normalsize{\@setfontsize\normalsize{9}{12}}\makeatother
\AtBeginDocument{\color{body}\normalsize}

\setlength{\parindent}{0pt}
\setlength{\parskip}{4pt}
\raggedcolumns
% Les listes posent une pénalité négative (−51) avant et après elles : avec des
% colonnes non étirées, c'est la coupure la moins chère, et multicols coupe là en
% laissant un vide. On la neutralise, et on interdit les lignes orphelines ou veuves
% (LaTeX recharge \clubpenalty depuis \@clubpenalty à chaque paragraphe).
\makeatletter
\@beginparpenalty=0 \@endparpenalty=0 \@itempenalty=0
\@clubpenalty=10000 \clubpenalty=10000 \widowpenalty=10000
\makeatother
\setlength{\columnsep}{8mm}
\arrayrulecolor{line}
\renewcommand{\tabularxcolumn}[1]{m{#1}}
\renewcommand{\arraystretch}{1.3}

\pagestyle{fancy}\fancyhf{}
\renewcommand{\headrulewidth}{0pt}
\renewcommand{\footrulewidth}{0pt}
\fancyfoot[L]{\scriptsize\color{muted}CY-HACK · Developer feedback}
\fancyfoot[R]{\scriptsize\color{muted}\thepage\,/\,\pageref{LastPage}}

% ── briques ──
% \leavevmode : un \color posé en mode vertical (item ou paragraphe qui commence par du
% code) laisse un whatsit dans la liste verticale, donc un point de coupure légal juste
% après un intitulé. En mode horizontal, le whatsit reste dans la ligne.
\newcommand{\hh}[1]{\par\addvspace{14pt}{\raggedright\hdfont\color{navy}\fontsize{10.5}{12.5}\selectfont\bfseries #1\par}\nobreak\addvspace{3pt}}
\newcommand{\hs}[1]{\leavevmode{\color{navy}\bfseries #1}}
\newcommand{\tm}[1]{\leavevmode{\ttfamily\color{navy}#1}}
\newcommand{\eyebrow}[1]{{\color{blue}\bfseries\fontsize{6.6}{8}\selectfont\addfontfeature{LetterSpace=6}#1}}

% liens : transaction dans l'explorer (hash abrégé à 8 caractères), fichier du repo
\newcommand{\explorer}{https://custom.xrpl.org/lending-hackathon.dev.ripplex.io:51233/transactions/}
\newcommand{\tx}[1]{\leavevmode\href{\explorer#1}{{\ttfamily\color{blue}\StrLeft{#1}{8}…}}}
\newcommand{\repo}[2]{\leavevmode\href{https://github.com/charlyppr/xrpl-lending/blob/main/#1}{{\ttfamily\color{blue}#2}}}

% sévérité : seule pastille et seule icône du document
\newcommand{\pill}[4]{\tcbox[on line,colback=#1,colframe=#1,boxrule=0pt,boxsep=0pt,left=3.2pt,right=3.2pt,
  top=1.5pt,bottom=1.5pt,arc=1.6pt]{\fontsize{6.4}{7}\selectfont\color{#2}#3\hspace{2.6pt}\bfseries #4}}
\newcommand{\sevHigh}{\pill{highbg}{highfg}{\faExclamationTriangle}{HIGH}}
\newcommand{\sevMedium}{\pill{medbg}{medfg}{\faExclamationCircle}{MEDIUM}}
\newcommand{\sevLow}{\pill{lowbg}{lowfg}{\faInfoCircle}{LOW}}

% en-tête d'un constat : numéro et titre, puis sévérité, catégorie, librairie
% #1 numéro · #2 High|Medium|Low · #3 catégorie · #4 titre
% after={\par\nobreak…} : le \nobreak suit directement la carte (avec after skip, le
% ressort posé entre la carte et le \nobreak laissait une coupure légale sous l'en-tête)
\newcommand{\finding}[4]{\par\addvspace{15pt}\noindent
  \IfStrEq{#2}{High}{\colorlet{sevrule}{highrule}}{\IfStrEq{#2}{Medium}{\colorlet{sevrule}{medrule}}{\colorlet{sevrule}{lowrule}}}%
  \begin{tcolorbox}[enhanced,colback=card,colframe=card,boxrule=0pt,arc=2pt,
    left=7pt,right=6pt,top=5pt,bottom=5.5pt,boxsep=0pt,before skip=0pt,
    after={\par\nobreak\nointerlineskip\vskip\glueexpr 5pt-\parskip\relax},
    borderline west={2.6pt}{0pt}{sevrule}]
  \hypertarget{f#1}{}\raggedright
  {\hangindent=13pt\hangafter=1\noindent\makebox[13pt][l]{\hdfont\color{blue}\fontsize{10}{12}\selectfont\bfseries #1}{\hdfont\color{navy}\fontsize{10}{12}\selectfont\bfseries #4}\par}
  \vspace{4pt}\hspace{13pt}\csname sev#2\endcsname\hspace{6pt}{\fontsize{7.2}{8}\selectfont\color{body}#3\hspace{4pt}·\hspace{4pt}\texttt{xrpl@4.6.0}}
  \end{tcolorbox}}

% intitulés en ligne, sans icône
\newcommand{\lbl}[1]{\hs{#1}\enspace}
\makeatletter
% \@afterheading : l'intitulé reste collé à la liste, comme un titre de section
\newcommand{\repro}{\par\addvspace{3pt}\lbl{Repro steps}\par\nobreak\vspace{0.5pt}\@afterheading}
\makeatother
\newcommand{\txs}{\par\addvspace{2pt}\lbl{Transactions}}
\newcommand{\fix}[1]{\par\addvspace{5pt}\begin{tcolorbox}[colback=bluecard,colframe=bluecard,boxrule=0pt,arc=2pt,
  left=6pt,right=6pt,top=4.5pt,bottom=4.5pt,boxsep=0pt,before skip=1pt,after skip=0pt]\raggedright
  \lbl{Proposed fix}#1\end{tcolorbox}\par}
\newenvironment{steps}{\begin{enumerate}[leftmargin=12pt,itemsep=1.5pt,topsep=2pt,parsep=0pt,beginpenalty=10000,
  label={\color{blue}\bfseries\fontsize{7.8}{8}\selectfont\arabic*}]}{\end{enumerate}}
\newcommand{\tabhead}[1]{{\bfseries\color{navy}#1}}

% ── captures d'écran de l'explorer ──
% Toutes au même grossissement : 460 px CSS (2415 px à ×5,25) remplissent la largeur
% utile de la colonne. Pas de titre de transaction : la pastille de statut seule.
\newtcolorbox{screen}{enhanced,colback=black,colframe=black,boxrule=0pt,arc=2.4pt,
  left=3.5pt,right=3.5pt,top=3pt,bottom=3pt,boxsep=0pt,before skip=0pt,after skip=2pt}
\newcommand{\shot}[1]{\includegraphics[scale=\fpeval{\linewidth/2415bp}]{#1.png}\par}
\newcommand{\capt}[1]{{\raggedright\fontsize{6.6}{8}\selectfont\color{muted}#1\par}}
% capture + légende : un seul bloc insécable, jamais coupé entre deux colonnes
\newenvironment{shotblock}{\par\addvspace{8pt}\noindent\begin{minipage}{\linewidth}\setlength{\parskip}{0pt}}{\end{minipage}\par\addvspace{2pt}}

\begin{document}

% ═══════════ BANDEAU ═══════════
\begin{tcolorbox}[colback=navy,colframe=navy,boxrule=0pt,arc=2.5pt,left=10pt,right=10pt,top=8pt,bottom=8pt,boxsep=0pt]
\begin{minipage}[c]{0.7\linewidth}
{\hyphenpenalty=10000\hdfont\color{white}\fontsize{15}{17}\selectfont\bfseries XRPL Lending Protocol · Developer Feedback}\\[3pt]
{\fontsize{7.4}{9}\selectfont\color{sky}DeVinci Blockchain × Ripple hackathon · 12 and 13 September 2026}
\end{minipage}\hfill
\begin{minipage}[c]{0.28\linewidth}\raggedleft
{\hdfont\color{white}\fontsize{9.5}{11}\selectfont\bfseries CY-HACK}\\[2pt]
{\fontsize{6.8}{8.4}\selectfont\color{sky}DevEx pseudonym \texttt{late-quail-92}\\
\href{https://github.com/charlyppr/xrpl-lending}{\texttt{github.com/charlyppr/xrpl-lending}}}
\end{minipage}
\end{tcolorbox}

% ═══════════ EN-TÊTE EXIGÉ PAR LA CONSIGNE ═══════════
\vspace{8pt}
{\setlength{\tabcolsep}{0pt}\fontsize{7.6}{9.6}\selectfont\renewcommand{\arraystretch}{1}\renewcommand{\tabularxcolumn}[1]{p{#1}}
\begin{tabularx}{\linewidth}{@{}>{\raggedright\arraybackslash}p{0.17\linewidth}>{\raggedright\arraybackslash}p{0.22\linewidth}>{\raggedright\arraybackslash}X>{\raggedright\arraybackslash}p{0.16\linewidth}@{}}
\eyebrow{TRACK} & \eyebrow{FLAVOUR} & \eyebrow{ENVIRONMENT} & \eyebrow{LIBRARY AND VERSION}\\[3pt]
\hs{Track 1}, open-ended vault &
\hs{Loaded}: XLS-65, XLS-66, TokenEscrow &
\hs{Custom Hackathon Devnet}, \tm{rippled} 3.4.0-rc1 (hackathon branch), \tm{LendingProtocolV1\_1} and \tm{fixCleanup3\_4\_0} enabled &
\hs{\texttt{xrpl@4.6.0}}\\
\end{tabularx}}

\vspace{8pt}
{\fontsize{7.4}{9.6}\selectfont\color{body}The 143 transactions cited are validated on the ledger, type and result code
included (\repo{bonus/scripts/_probe-audit-hashes.mjs}{\_probe-audit-hashes.mjs}). Documentation claims are checked against XLS-65, XLS-66 and
xrpl.org, library claims against the \tm{rippled} and \tm{xrpl.js} source. Hashes link to the explorer.\par}

% ═══════════ SYNTHÈSE ═══════════
\vspace{8pt}
{\renewcommand{\arraystretch}{1.42}\setlength{\tabcolsep}{4pt}\fontsize{7.8}{9.4}\selectfont
\begin{tabularx}{\linewidth}{@{}r >{\raggedright\arraybackslash}X l l@{}}
\toprule
\tabhead{\#} & \tabhead{Finding} & \tabhead{Category} & \tabhead{Severity}\\
\midrule
\hyperlink{f1}{\color{blue}\bfseries 1} & \hyperlink{f1}{\tm{fixCleanup3\_4\_0} is enabled and missing from xrpl.org} & client libraries & \sevHigh\\
\hyperlink{f2}{\color{blue}\bfseries 2} & \hyperlink{f2}{A late \tm{LoanPay} without \tm{tfLoanLatePayment} returns \tm{tecEXPIRED}} & documentation/tutorials & \sevHigh\\
\hyperlink{f3}{\color{blue}\bfseries 3} & \hyperlink{f3}{First-loss capital covered 0.5\,\% of a defaulted loan} & UX & \sevHigh\\
\hyperlink{f4}{\color{blue}\bfseries 4} & \hyperlink{f4}{xrpl.org omits the sole-holder exception, and zero-valued fields are absent} & documentation/tutorials & \sevMedium\\
\hyperlink{f5}{\color{blue}\bfseries 5} & \hyperlink{f5}{No read command for loans or brokers} & missing primitive & \sevMedium\\
\hyperlink{f6}{\color{blue}\bfseries 6} & \hyperlink{f6}{One result code covers causes that need different actions} & UX & \sevMedium\\
\hyperlink{f7}{\color{blue}\bfseries 7} & \hyperlink{f7}{The hackathon build reverts the V1.1 \tm{LoanBrokerSet} restriction} & documentation/tutorials & \sevMedium\\
\hyperlink{f8}{\color{blue}\bfseries 8} & \hyperlink{f8}{Vault shares carry \tm{lsfMPTCanTrade}, but \tm{OfferCreate} rejects them} & documentation/tutorials & \sevMedium\\
\hyperlink{f9}{\color{blue}\bfseries 9} & \hyperlink{f9}{No flag closes a vault to deposits} & missing primitive & \sevLow\\
\bottomrule
\end{tabularx}}
\vspace{5pt}
{\fontsize{6.9}{8.8}\selectfont\color{muted}\hs{High}: blocks a minimum bar step, or misstates a depositor's risk. \hs{Medium}: costs time or a wrong design, with a workaround. \hs{Low}: documented elsewhere, or cosmetic.\par}

\vspace{4pt}
\begin{multicols}{2}\RaggedRight

\hh{The brief's questions}
{\renewcommand{\arraystretch}{1.35}\setlength{\tabcolsep}{0pt}\fontsize{8}{10.2}\selectfont\renewcommand{\tabularxcolumn}[1]{p{#1}}
\begin{tabularx}{\linewidth}{@{}>{\raggedright\arraybackslash}p{0.34\linewidth}>{\raggedright\arraybackslash}X@{}}
\hs{Vault, broker, loan} & Clear. Only the vault owner can submit \tm{LoanBrokerSet}. The drawdown has no transaction: \tm{LoanSet} moves the principal (\tx{D19CD59FB93CCF28547533C76CC1AFA71CF592FD0EC9D0A9E217A9F6EAAD85DA}).\\
\hs{First-loss parameters} & Not as named (\hyperlink{f3}{3}).\\
\hs{Multi-party \tm{LoanSet}} & Broker signs, borrower countersigns and submits. The \tm{xrpl@4.6.0} helper fails here (\hyperlink{f1}{1}).\\
\hs{SDK or raw JSON} & All 15 XLS-65/66 types are typed. No raw JSON.\\
\hs{Position and yield} & Computed client-side from five calls (\hyperlink{f5}{5}).\\
\hs{Docs and explorer} & The docs differ in \hyperlink{f1}{1}, \hyperlink{f2}{2}, \hyperlink{f4}{4}, \hyperlink{f8}{8}. The explorer hides a flag name (\hyperlink{f2}{2}) and \tm{Vault} fields (\hyperlink{f4}{4}).\\
\end{tabularx}}

% ───────────── 1 ─────────────
\finding{1}{High}{client libraries}{\tm{fixCleanup3\_4\_0} is enabled and missing from xrpl.org}

The amendment is active on the devnet and absent from \emph{Known Amendments}. It changes two
lending behaviours, and neither the tutorial nor \tm{xrpl@4.6.0} reflects them.

\hs{Impairment.} XLS-66, section 3.10.4.2, condition 9 rejects \tm{tfLoanImpair} while
\tm{currentTime <= NextPaymentDueDate}. The tutorial \emph{Manage a Loan} still says a broker
can impair ``\emph{before a payment due date passes}''.

\hs{Counterparty signature.} Since rippled PR \#8162, \tm{CounterpartySignature} is signed with
the \tm{CPT} prefix. \tm{signLoanSetByCounterparty} signs with \tm{encodeForSigning}, the
previous prefix, and offers no option.

\repro
\begin{steps}
\item \tm{LoanSet} with \tm{PaymentInterval} 60 and \tm{GracePeriod} 60. \tm{LoanManage} \tm{tfLoanImpair} returns \tm{tecTOO\_SOON} 31 s and 13 s before \tm{NextPaymentDueDate}, and \tm{tesSUCCESS} 5 s after, by the script's clock. The ledger closes the accepted transaction 8 s after the due date.
\item The broker signs a \tm{LoanSet}, the borrower applies \tm{signLoanSetByCounterparty} and submits: \tm{Counterparty: Invalid signature}, a local check with no hash.
\item The same blob signed over \tm{encodeForSigningCounterparty} (\repo{scripts/raw-submit.mjs}{raw-submit.mjs}) returns \tm{tesSUCCESS}.
\end{steps}

\txs impairment \tx{5486020BA8FF227C58409042F96BDC07246FDACC3F9B0EF24749555E319E972E}, countersigned \tm{LoanSet} \tx{23631C3610DA97926E8BF9FC5BC4E76E504C9DFB2683C928E825DE0138397B2B}

\fix{List \tm{fixCleanup3\_4\_0} on \emph{Known Amendments} with both lending changes. Update \emph{Manage a Loan}: impairment is accepted once \tm{NextPaymentDueDate} has passed. In \tm{xrpl.js}, sign \tm{CounterpartySignature} with \tm{encodeForSigningCounterparty}, keeping the previous prefix as an option for networks without the amendment.}

% ───────────── 2 ─────────────
\finding{2}{High}{documentation/tutorials}{A late \tm{LoanPay} without \tm{tfLoanLatePayment} returns \tm{tecEXPIRED}}

Once \tm{NextPaymentDueDate} has passed, every \tm{LoanPay} without the flag fails, even within
the grace period: the instalment, the instalment with \tm{LatePaymentFee}, the full balance, and
\tm{tfLoanFullPayment}. XLS-66, section 3.11.4.2, condition 11 defines this failure; the xrpl.org
\tm{LoanPay} reference lists eight error codes, not this one. The explorer shows the flag as
\tm{0x00040000}, and no \tm{Loan} field holds the accrued late interest.

\repro
\begin{steps}
\item \tm{LoanSet} with \tm{PaymentInterval} 90, \tm{GracePeriod} 60, \tm{PaymentTotal} 4 (\repo{bonus/scripts/_probe-pay-states.mjs}{\_probe-pay-states.mjs}).
\item 11 s after the due date by ledger close time, \tm{LoanPay} of \tm{ceil(PeriodicPayment) + LoanServiceFee + LatePaymentFee} returns \tm{tecEXPIRED} with \tm{Flags: 0}, and \tm{tesSUCCESS} with \tm{Flags: 262144}.
\item With \tm{LateInterestRate} 30\,\% (\repo{bonus/scripts/_probe-h2-race.mjs}{\_probe-h2-race.mjs}), the same amount returns \tm{tecINSUFFICIENT\_PAYMENT}. Sending 2\,765\,001 drops succeeds and takes 765\,003.
\end{steps}

\txs no flag \tx{96C3870480E4CB0C9E3B925F2965347DA0815A486CDB4B49545928E32D6B0BA4}, flag \tx{EFAD383F9B622357CE67CBB0518D9F9F5FC27BE611D43F26FD69427713FB608F}, late interest \tx{9A7589384C5142703AFFACBB91B12E223EC734F7ACBE51272A87525FE4EE9CA8}

\fix{Add \tm{tecEXPIRED} to the \tm{LoanPay} error table, with \tm{tfLoanLatePayment} as the remedy. Name the flag in the explorer. Expose the amount due now, late interest included, in a \tm{Loan} field or a read command.}

\begin{shotblock}
\begin{screen}
\shot{f1a-badge}\shot{f1a-status}\vspace{3pt}\shot{f1b-badge}\shot{f1b-flags}
\end{screen}
\capt{Explorer. Same loan and amount, without the flag, then with it.}
\end{shotblock}

% ───────────── 3 ─────────────
\finding{3}{High}{UX}{First-loss capital covered 0.5\,\% of a defaulted loan}

Two brokers, both with \tm{CoverRateMinimum} 10\,\% and \tm{CoverRateLiquidation} 5\,\%:

{\setlength{\tabcolsep}{3pt}\fontsize{7.4}{9}\selectfont
\begin{tabularx}{\linewidth}{@{}r r r r r >{\raggedleft\arraybackslash}X@{}}
\toprule
\tabhead{Vault} & \tabhead{Loan} & \tabhead{Cover} & \tabhead{Taken} & \tabhead{Loss} & \tabhead{Share value}\\
\midrule
200 XRP & 50 & 50 & 0.25 & 49.75 & 1.000\,$\rightarrow$\,0.751\\
5 XRP & 4 & 1 & 0.02 & 3.98 & 1.000\,$\rightarrow$\,0.204\\
\bottomrule
\end{tabularx}}

The cover taken is \tm{debt $\times$ 10\% $\times$ 5\%}, to the drop. This is the documented
formula: \tm{CoverRateLiquidation} applies to the minimum cover, not to the defaulted debt, and
the xrpl.org example covers 1\,\% of the debt. After the default \tm{DebtTotal} is zero, so is
the cover floor, and the broker withdrew the rest. Second vault, UTC:\par\nobreak

{\setlength{\tabcolsep}{3pt}\fontsize{7.4}{9}\selectfont
\begin{tabularx}{\linewidth}{@{}l >{\raggedright\arraybackslash}X r@{}}
\toprule
14:23:50 & \tm{VaultWithdraw}: \tm{tecINSUFFICIENT\_FUNDS} & \tx{D91BA47E816FDA60F9BBAD7F100B383099D87021B08DE065E841B9E8E34AFE2D}\\
14:27:01 & \tm{tfLoanDefault}: 0.02 XRP from the cover & \tx{923F7D616B094C82197506CF2867907A93067041E3990D24BD07A19035F73EAC}\\
14:27:11 & \tm{VaultWithdraw}: 1.02 XRP & \tx{3338BA629F8772CDEC07B6F714F2AF3A1ADC8E93642429A54B95CF372FDB2880}\\
14:27:22 & \tm{LoanBrokerCoverWithdraw}: 0.98 XRP & \tx{68DD97D949602FD470E0FA80B2526C3FDCFA3449095F7A7879F4D4430845D45E}\\
\bottomrule
\end{tabularx}}

\repro
\begin{steps}
\item \repo{bonus/scripts/h1-default.mjs}{h1-default.mjs}, after \tm{setup-accounts.mjs}: a 200 XRP vault, 50 XRP of cover, one 50 XRP loan, defaulted after \tm{GracePeriod}.
\item Compare \tm{CoverAvailable} and \tm{AssetsTotal} before and after the default.
\item \tm{LoanBrokerCoverWithdraw} of the remaining cover returns \tm{tesSUCCESS}.
\end{steps}

\txs first vault default \tx{91F458E3E6244E1C3005345C5C83CDFAF0E8A1AF6E5587E100C302CA56743BC7}

\fix{Next to the formula, state the share of a defaulted debt the cover absorbs: \tm{CoverRateMinimum $\times$ CoverRateLiquidation}, 0.5\,\% at 10\,\% and 5\,\%. Hold \tm{LoanBrokerCoverWithdraw} for a set delay after a default, or state that no delay exists.}

% ───────────── 4 ─────────────
\finding{4}{Medium}{documentation/tutorials}{xrpl.org omits the sole-holder exception, and zero-valued fields are absent}

XLS-65 redeems shares at \tm{(AssetsTotal − LossUnrealized) / total shares}, except for the sole
outstanding holder, who is not charged \tm{LossUnrealized}. The xrpl.org Single Asset Vault page
gives the formula without the exception. \tm{LossUnrealized} is absent until an impairment
writes it, as are \tm{AssetsAvailable}, \tm{AssetsMaximum}, \tm{PaymentRemaining} and
\tm{TotalValueOutstanding} at zero. The explorer's \emph{Detailed} tab shows no field of the
\tm{Vault} node.

\repro
\begin{steps}
\item \tm{vault\_info} on a vault without an impaired loan: no \tm{LossUnrealized} field.
\item \tm{tfLoanImpair} on one of its loans, after \tm{NextPaymentDueDate}.
\item \tm{vault\_info} returns \tm{LossUnrealized} \tm{"3000001"} with \tm{AssetsTotal} \tm{"5000002"}. Compare the explorer's \emph{Detailed} and \emph{Raw} tabs.
\end{steps}

\txs impairment \tx{C8678C1C0A10BE889124F40C03017038654B278E2D464FAC54784F55B35318FC}

\fix{Add the sole-holder exception to the xrpl.org formula. Return zero-valued fields as \tm{"0"} in \tm{vault\_info}, or document their absence. Show changed field values in the explorer's \emph{Detailed} tab.}

\begin{shotblock}
\begin{screen}
\shot{f2-badge}\shot{f2-meta}\vspace{3pt}\shot{f2-raw}
\end{screen}
\capt{Explorer. The impairment: \emph{Detailed} tab, then \emph{Raw} tab with the \texttt{Vault} node expanded.}
\end{shotblock}

% ───────────── 5 ─────────────
\finding{5}{Medium}{missing primitive}{No read command for loans or brokers}

A depositor's position takes five calls, then six values computed client-side in
\repo{scripts/lib/nav.mjs}{lib/nav.mjs}: gross and net share value, position value, withdrawable
amount, utilisation, broker debt capacity. \tm{mpt\_holders} is documented as Clio-only, and the
devnet serves \tm{rippled}.\par\nobreak

{\setlength{\tabcolsep}{3pt}\fontsize{7.4}{9}\selectfont
\begin{tabularx}{\linewidth}{@{}l >{\raggedright\arraybackslash}X@{}}
\toprule
\tabhead{Call} & \tabhead{Returns}\\
\midrule
\tm{vault\_info} & totals, share issuance and flags\\
\tm{account\_objects} \tm{mptoken} & the depositor's shares\\
\tm{account\_objects}, owner & the broker, filtered on \tm{VaultID}\\
\tm{account\_objects} \tm{loan} & the loans, on the broker's pseudo-account\\
\tm{mpt\_holders} & \tm{unknownCmd}\\
\bottomrule
\end{tabularx}}

\repro
\begin{steps}
\item Send \tm{loan\_info} and \tm{loan\_broker\_info} to the devnet JSON-RPC endpoint: \tm{unknownCmd}. Neither has a documentation page.
\end{steps}

\fix{Add \tm{loan\_broker\_info} and \tm{loan\_info}, or list a vault's brokers and loans in \tm{vault\_info}. Return share value, withdrawable amount and utilisation. Expose Clio on the devnet.}

% ───────────── 6 ─────────────
\finding{6}{Medium}{UX}{One result code covers causes that need different actions}

A broker who receives \tm{tecINSUFFICIENT\_FUNDS} on \tm{LoanSet} cannot tell whether to add
cover or wait for deposits. Three other codes share the problem.\par\nobreak

{\setlength{\tabcolsep}{3pt}\fontsize{7.2}{8.8}\selectfont
\begin{tabularx}{\linewidth}{@{}l >{\raggedright\arraybackslash}X r@{}}
\toprule
\tabhead{Code} & \tabhead{Cause} & \\
\midrule
\tm{tecINSUFFICIENT\_FUNDS} & \tm{LoanSet} above the cover limit & \tx{5DC3A6E00F31DA82C69AB77669C64691BD71CAC49C169736F53D7C724147618B}\\
 & \tm{LoanSet} above the vault liquidity & \tx{B92A90F59EBF79108B5414F61A12656F321EBC264E1BB186FCC14CAF259F9CAA}\\
 & \tm{LoanBrokerCoverWithdraw} below the floor & \tx{4FB4A2B0D01DD2411A5E8592680D50B1F4CCA51ED97C243C2A506516262F1370}\\
\midrule
\tm{tecLIMIT\_EXCEEDED} & \tm{VaultDeposit} above \tm{AssetsMaximum} & \tx{132BAF7C1B8208E695A453C0EE061BC5E3528FC97571589D8A2EA66F027B5B4F}\\
 & \tm{AssetsMaximum} set below \tm{AssetsTotal} & \tx{37053F15A20E6153EDBAD2DE6DD48E039E7A234328382A58DD97607FA88C3DC6}\\
 & \tm{DebtMaximum} set below \tm{DebtTotal} & \tx{82C14ED4071391EC1EF25637D2FEC28CF784BB19FEB4E89082A64CD9748E4121}\\
\midrule
\tm{tecNO\_PERMISSION} & \tm{tfLoanImpair} sent by the borrower & \tx{F030419492FB6A9A8CD4BA3015E8A55A9A86B9A1288A6458C5CC36475DE2C112}\\
 & \tm{tfLoanUnimpair} on an unimpaired loan & \tx{13975FD33EFA2F97D3120920E6F613BAFB95D09459350E7693C207E66AA24537}\\
\midrule
\tm{tecNO\_AUTH} & shares sent to a holder without \tm{MPToken} & \tx{789AAF7E228CC085199E2FB38B07A10DFC68A45C76C5FC3EFDB498BC5FD512A4}\\
 & non-transferable shares sent & \tx{E9DE504D40B8AB27F7C7EA433B0C1E68DC91854AA04A5D600163A64DB8BE5FD9}\\
\bottomrule
\end{tabularx}}

\repro
\begin{steps}
\item \repo{bonus/scripts/_probe-cliffs.mjs}{\_probe-cliffs.mjs}: a \tm{LoanSet} above the cover limit with ample liquidity, then one above the liquidity with ample cover. Both return \tm{tecINSUFFICIENT\_FUNDS}.
\end{steps}

\fix{Where one transaction type maps two causes to one code (\tm{LoanSet}, \tm{LoanManage}, a \tm{Payment} of shares), return distinct codes or a result message naming the check that failed.}

% ───────────── 7 ─────────────
\finding{7}{Medium}{documentation/tutorials}{The hackathon build reverts the V1.1 \tm{LoanBrokerSet} restriction}

The V1.1 documentation, \href{https://opensource.ripple.com/docs/lending-protocol-v1-1/closed-ended-vaults}{\emph{Closed-ended vaults}} on opensource.ripple.com, restricts \tm{LoanBrokerSet} on open-ended vaults, a rule added by
rippled PR \#8076. A revert of \#8076 reached the lending hackathon branch on 10 September 2026.
The other V1.1 rules tested (new fields, date validation, phase locks) are enforced. The event
brief does not mention the revert.

\repro
\begin{steps}
\item \tm{feature}: \tm{LendingProtocolV1\_1} is enabled.
\item \tm{LoanBrokerSet} on an open-ended vault returns \tm{tesSUCCESS}.
\item \tm{LoanBrokerSet} on a \tm{VaultKind: 1} vault returns \tm{tesSUCCESS} (\repo{bonus/scripts/_probe-h13-closed.mjs}{\_probe-h13-closed.mjs}).
\end{steps}

\txs open-ended \tx{781F54B5E77A768F4C02A54E3C55902AA9F1AC96AA04037AB97CE93FFB5DBDE4}, closed-ended \tx{1961BE21CC4011F50AA926E1494E9D51F94FC57B0E9C0C1B8F8ADF6FEDD09D40}

\fix{State in the event brief which V1.1 rules the hackathon build enforces, and that the \#8076 restriction is reverted.}

% ───────────── 8 ─────────────
\finding{8}{Medium}{documentation/tutorials}{Vault shares carry \tm{lsfMPTCanTrade}, but \tm{OfferCreate} rejects them}

\tm{VaultCreate} sets \tm{lsfMPTCanEscrow}, \tm{lsfMPTCanTrade} and \tm{lsfMPTCanTransfer} on the
share issuance, with no field to change them. The xrpl.org \tm{MPTokenIssuance} reference says
\tm{lsfMPTCanTrade} lets holders trade ``\emph{using the XRP Ledger DEX or AMM}''. No amendment
enabled on the devnet opens the DEX to MPTs.

\repro
\begin{steps}
\item \tm{vault\_info}: \tm{shares.Flags} is 56, that is 0x08, 0x10 and 0x20.
\item \tm{OfferCreate} by a holder, \tm{TakerGets} the share MPT, \tm{TakerPays} XRP: \tm{temDISABLED}, a local rejection with no hash (\repo{bonus/scripts/_probe-shares.mjs}{\_probe-shares.mjs}).
\end{steps}

\fix{State on the \tm{MPTokenIssuance} reference that DEX trading of MPTs needs an amendment not yet available, or leave \tm{lsfMPTCanTrade} unset on vault shares until then.}

% ───────────── 9 ─────────────
\finding{9}{Low}{missing primitive}{No flag closes a vault to deposits}

\tm{VaultSet} exposes \tm{Data}, \tm{AssetsMaximum} and \tm{DomainID}. The only way to stop
deposits is a cap equal to \tm{AssetsTotal}. \tm{AssetsMaximum: 0} means no cap in XLS-65 and
the \tm{Vault} entry; the \tm{VaultSet} reference does not repeat it.

\repro
\begin{steps}
\item \tm{VaultSet} with \tm{AssetsMaximum} equal to \tm{AssetsTotal} returns \tm{tesSUCCESS}.
\item A 1 XRP \tm{VaultDeposit} returns \tm{tecLIMIT\_EXCEEDED}.
\end{steps}

\txs deposit refused \tx{132BAF7C1B8208E695A453C0EE061BC5E3528FC97571589D8A2EA66F027B5B4F}

\fix{Add a \tm{VaultSet} flag that closes a vault to deposits. Repeat that \tm{0} means no cap on the \tm{VaultSet} reference.}

% ───────────── fin ─────────────
\hh{Minor issues}
{\renewcommand{\arraystretch}{1.35}\setlength{\tabcolsep}{3pt}\fontsize{7.4}{9}\selectfont
\begin{tabularx}{\linewidth}{@{}>{\raggedright\arraybackslash}X l l@{}}
\toprule
\tabhead{Issue} & \tabhead{Category} & \tabhead{Severity}\\
\midrule
The Lending Protocol concept page spells \tm{PrincipleOutstanding} and \tm{depostitor}. & documentation/tutorials & \sevLow\\
\tm{validateVaultCreate} in \tm{xrpl@4.6.0} accepts any \tm{WithdrawalPolicy}; 0, 2, 3, 99 and 255 return \tm{temMALFORMED}. & client libraries & \sevLow\\
The faucet ignores \tm{destination} and funds a new account instead. & other & \sevLow\\
\bottomrule
\end{tabularx}}

\hh{What worked}
\tm{tecHAS\_OBLIGATIONS} names its cause on the three delete transactions. \tm{VaultDelete}
removes third-party \tm{MPToken} objects and refunds their reserve. The cover floor is enforced
to the drop, and cover rates cannot change on an existing broker (\tm{temINVALID}). Closed-ended
vault phases behave as XLS-65 describes.

\hs{Limits.} XRP only, so no clawback with an issued asset. Payment intervals of 60 to 90 s. One
build. The minimum bar run, the guardrails and the escrow swap are in the
\href{https://github.com/charlyppr/xrpl-lending\#readme}{\color{blue}README}.

\end{multicols}
\end{document}
